\chapter{Alpha-1 Antitrypsin}
\section*{What Is This Test?}
Alpha-1 antitrypsin (AAT) is a 52-kDa glycoprotein. It is synthesized primarily by hepatocytes. It is also made, to a lesser extent, by alveolar macrophages and intestinal epithelial cells. It is the most abundant serine protease inhibitor (serpin) in plasma. Its main role is to inhibit neutrophil elastase. Neutrophil elastase is a proteolytic enzyme released by activated neutrophils. It degrades elastin in the alveolar walls. AAT deficiency is an autosomal codominant genetic disorder. It is caused by mutations in the SERPINA1 gene (chromosome 14q32.1). The most common pathogenic variant is the Z allele (Glu342Lys, p.Glu366Lys in the mature protein). It causes polymerization of the AAT protein within hepatocyte endoplasmic reticulum. This has two results. (1) Secretion of AAT into the blood is reduced, producing low serum AAT levels. This leads to unopposed neutrophil elastase activity in the lungs. It causes progressive destruction of alveolar elastin and early-onset panacinar emphysema. The risk is especially high in smokers, who have increased neutrophil elastase burden from chronic inflammation. Basal predominance is seen on CT. (2) Polymerized Z-AAT accumulates in hepatocyte ER. This causes chronic hepatocellular injury. It leads to neonatal cholestasis, chronic hepatitis, cirrhosis, and hepatocellular carcinoma (HCC) in adulthood. The normal (wild-type) allele is the M allele. The most common genotypes are Pi*MM, Pi*MZ, Pi*SZ, Pi*ZZ, and Null/Null. Pi*MM has normal AAT levels and accounts for >90\% of the population. Pi*MZ is a heterozygous carrier with AAT ~60\% of normal. It carries a slightly increased risk of lung disease in smokers. Pi*SZ is a compound heterozygote with AAT ~40\%. Pi*ZZ is a homozygous Z with AAT ~15\% of normal. It is a severe AAT deficiency with high risk of emphysema and liver disease. Null/Null produces no AAT. It causes severe emphysema without liver disease, with no protein accumulation in the ER.
\section*{Alternative Names}
\begin{itemize}
  \item AAT.
  \item Alpha-1 Antitrypsin Level.
  \item AAT Phenotype.
  \item A1AT.
  \item SERPINA1.
  \item Pi Typing.
\end{itemize}
\section*{Principle}
Serum AAT level is measured by immunoturbidimetry or immunonephelometry using anti-AAT antibodies. AAT phenotype (Pi typing) is determined by isoelectric focusing (IEF) of serum proteins. IEF separates AAT variants by their isoelectric point (M, S, Z, and other rare variants). AAT genotyping (SERPINA1 gene sequencing) is the most definitive method for identifying pathogenic alleles. Complete diagnosis requires combining the serum AAT level with the AAT phenotype or genotype. Low AAT level with Pi*ZZ phenotype or genotype means severe AAT deficiency.
\section*{History}
Alpha-1 antitrypsin deficiency was discovered in 1963. The discoverers were the Swedish physician Carl-Bertil Laurell and his trainee Sten Eriksson. They noted the absence of the alpha-1 globulin band on serum protein electrophoresis in patients with emphysema \cite{laurell1963electrophoretic}. In 1965, Eriksson established the hereditary nature of the deficiency in a series of Swedish families with early-onset emphysema. In 1969, Sharp and colleagues reported that AAT deficiency also caused cirrhosis and neonatal liver disease \cite{sharp1969cirrhosis}. The molecular basis was clarified in subsequent decades. The SERPINA1 gene was mapped to chromosome 14. In 1992, Lomas and colleagues demonstrated that the Z variant polymerizes within hepatocytes. This explained the intracellular accumulation and liver injury \cite{lomas1992mechanism}. Intravenous augmentation therapy with plasma-derived AAT was approved in the United States in 1987. Targeted and newborn screening programs now permit early detection of the disorder.
\section*{Why This Test Is Ordered}
\begin{itemize}
  \item Evaluation of early-onset emphysema (before age 45), especially with basilar (panacinar) predominance or a minimal smoking history.
  \item Workup of unexplained chronic liver disease, cirrhosis, or neonatal cholestasis.
  \item Investigation of poorly reversible asthma, bronchiectasis, or chronic cough of unclear cause.
  \item Evaluation of necrotizing panniculitis or ANCA-positive vasculitis (granulomatosis with polyangiitis).
  \item Screening of first-degree relatives of patients with known AAT deficiency.
  \item Confirmatory testing when serum protein electrophoresis shows an absent or reduced alpha-1 globulin band.
\end{itemize}
\section*{Primary vs Secondary Test}
This is a diagnostic and confirmatory test rather than a routine screening test. It is ordered when clinical findings raise suspicion of AAT deficiency. These findings include early-onset emphysema, unexplained liver disease, and panniculitis. A complete diagnosis requires combining the serum AAT level with phenotype or genotype testing. Targeted screening of at-risk relatives of affected patients is recommended. Population-wide screening is not.
\section*{How This Test Is Performed}
A blood sample is collected by standard venipuncture into a serum separator (gold-top) tube for the quantitative AAT level. Serum AAT is measured by immunoturbidimetry or immunonephelometry. AAT phenotype (Pi typing) is determined by isoelectric focusing. Genotype is determined by PCR-based assays or SERPINA1 gene sequencing. Genotype testing typically uses a separate EDTA (purple-top) tube. Quantitative results are usually available within 1--2 days. Phenotype and genotype results take longer.
\section*{What Preparation Is Needed}
No fasting or special preparation is required. AAT is an acute phase reactant. Testing during an acute inflammatory illness or pregnancy can transiently raise the serum level. This can mask a mild deficiency. Results should be interpreted with this in mind. Patients should continue their usual medications unless directed otherwise.
\section*{Contraindications and Precautions}
\begin{itemize}
  \item There are no absolute contraindications to venipuncture. Interpretive cautions: a serum AAT level within the normal range does not exclude mild-to-moderate deficiency (Pi*MZ or Pi*SZ) when measured during an acute phase response. Phenotype or genotype testing is needed if clinical suspicion persists. Some rare null and dysfunctional alleles are not fully characterized by routine genotyping. Phenotype testing may be required. In patients receiving augmentation therapy, the measured serum AAT level reflects the infused protein. It should not be used to judge the diagnosis.
\end{itemize}
\section*{Risks}
Minimal risk. Slight pain or bruising at the venipuncture site. Rarely: hematoma, infection, or vasovagal syncope.
\section*{What Happens After the Test}
Results return within days and are interpreted by combining the serum AAT level with the phenotype or genotype. A confirmed diagnosis prompts evaluation of lung function (spirometry with diffusing capacity), chest imaging, and liver studies. It also prompts referral to a pulmonologist or hepatologist and genetic counseling for the family. Aggressive smoking cessation is essential. Smoking markedly accelerates the development of emphysema. In patients with established airflow obstruction from severe deficiency, intravenous augmentation therapy may be considered.
\section*{Understanding Results}
\subsection*{Low/Normal AAT}
Pi*MM: AAT 100--200 mg/dL (normal). Pi*MZ: AAT ~60\% of normal; borderline. Pi*SZ: AAT ~40\%; moderate deficiency, risk of emphysema in smokers.
\subsection*{Severely Low AAT}
Pi*ZZ: AAT <50--60 mg/dL (typically 15--30 mg/dL, ~15\% of normal). This is severe AAT deficiency with high risk of several conditions. These include panacinar emphysema, liver disease, necrotizing panniculitis (rare), and ANCA-positive vasculitis (granulomatosis with polyangiitis, GPA). Panacinar emphysema has onset age 30--50. It is markedly accelerated by smoking. Average life expectancy of smokers with Pi*ZZ is 50--55 years, vs. 65--70 years for non-smokers. Liver disease includes neonatal cholestatic jaundice, chronic hepatitis, cirrhosis, and HCC. Null/Null: AAT undetectable or near-zero. This causes severe emphysema without liver disease.
\section*{Normal Results}
Serum AAT: 100--200 mg/dL (immunoturbidimetry). AAT >57 mg/dL (11 mcmol/L) is the protective threshold. Below this threshold the risk of emphysema increases significantly. Pi*MM phenotype: normal AAT. AAT is an acute phase reactant. It rises 2--4 fold during inflammation, infection, pregnancy, and estrogen therapy. A normal AAT level during acute inflammation may mask mild-to-moderate AAT deficiency (Pi*MZ or Pi*SZ). Phenotype or genotype testing is definitive.
\section*{Follow-up Tests}
Follow-up tests include AAT phenotype (Pi typing by isoelectric focusing), SERPINA1 gene sequencing, pulmonary function tests (spirometry with DLCO), chest CT, liver function tests and ultrasound. Liver biopsy is used if hepatocellular injury is suspected.
\section*{References}
\bibliographystyle{unsrtnat}
\bibliography{references}

\clearpage
\section*{Regulatory Status and Authorization Date}
This chapter describes a test category, not a specific commercial product. FDA, CDSCO, EU IVDR, and MHRA approval or registration dates therefore cannot be assigned without the assay manufacturer, product name, instrument, and intended use. For laboratory-developed tests, an individual agency approval date may not exist. See the Preface for the interpretation of regulatory status.

\clearpage
